\chapter{Аналитическая часть}

В данном разделе будут представлены описания алгоритмов нахождения расстояний Левенштейна и Дамерау~---~Левенштейна и их практическое применение.

\section{Расстояние Левенштейна}

Расстояние Левенштейна~--- это метрика между двумя строками. Оно равно минимальному числу элементарных операций (вставок, удалений и замен символов), необходимых для преобразования одной строки в другую\cite{lev-damlev-analysis}.
Каждая операция имеет свою цену($w$). Редакционное предписание~--- это последовательность операций с суммарной минимальной стоимостью, которую необходимо выполнить для получения из первой строки вторую. Эта цена и есть искомое расстояние Левенштейна.

Используются следующие обозначения:
\begin{enumerate}
    \item I~(от англ. insert)~--- вставка ($w(\lambda, b) = 1$);
    \item R~(от англ. replace)~--- замена ($w(a, b) = 1$, $a \neq b$);
    \item D~(от англ. delete)~--- удаление ($w(a, \lambda) = 1$).
\end{enumerate}

Функция $D(i, j)$ определяет редакционное расстояние между строками $S_1[1...i]$ и $S_2[1...j]$.

Расстояние Левенштейна между двумя строками $S_{1}$ и $S_{2}$ (длиной $M$ и $N$ соответственно) рассчитывается по следующей рекуррентной формуле:
\begin{equation}
	\label{eq:L}
	D(i, j) =
	\begin{cases}
		0, &\text{i = 0, j = 0}\\
		i, &\text{j = 0, i > 0}\\
		j, &\text{i = 0, j > 0}\\
		min \begin{cases}
			D(i, j - 1) + 1,\\
			D(i - 1, j) + 1,\\
			D(i - 1, j - 1) +  m(S_{1}[i], S_{2}[j]), \\
		\end{cases}
		&\text{i > 0, j > 0}
	\end{cases}
\end{equation}
где сравнение символов строк $S_1$ и $S_2$ рассчитывается таким образом:
\begin{equation}
	\label{eq:m}
	m(a, b) = \begin{cases}
		0 &\text{если a = b,}\\
		1 &\text{иначе.}
	\end{cases}
\end{equation}

\section{Нерекурсивный алгоритм нахождения расстояния Левенштейна}

Рекурсивная реализация алгоритма нахождения расстояния Левенштейна малоэффективна по времени при больших $N, M > 10$, так как производится много повторных вычислений.
Для оптимизации применяется матрица размерности $(N + 1) \times (M + 1)$, в которой будут храниться промежуточные значения. Первый элемент матрицы заполен нулём.
Остальные требуется заполнить в соответствии с формулой (\ref{eq:L}). Результатом работы будет являться значение в последней ячейке последний строки матрицы.
Для вычислений необязательно хранить всю матрицу --- достаточно двух строк.


\section{Рекурсивный алгоритм нахождения расстояния Левенштейна}

Рекурсивный алгоритм реализует формулу (\ref{eq:L}).

\section{Рекурсивный алгоритм нахождения расстояния Левенштейна с кэшем}

Рекурсивный алгоритм заполнения оптимизируется по времени выполнения с использованием кеша. В качестве кеша используется матрица.
В случае, если рекурсивный алгоритм выполняет расчёт для данных, которые ещё не были обработаны (в ячейке находится бесконечность), результат нахождения расстояния заносится в матрицу.
В случае, если ячейка не равна бесконечности, расстояние не находится и алгоритм переходит к следующему шагу.

\section{Расстояние Дамерау~---~Левенштейна}

Расстояние Дамерау~---~Левенштейна~--- это метрика между двумя строками. Оно равно минимальному числу элементарных операций (вставок, удалений замен и перестановок символов), необходимых для преобразования одной строки в другую~\cite{lev-damlev-analysis}. Является расширением расстояния Левенштейна, поскольку помимо трёх базовых операций содержит ещё операцию транспозиции $T$ (от англ. transposition).

Расстояние Дамерау~---~Левенштейна определятся следующей рекуррентной формуле:

\begin{multline}
    D(m, n) =\\ =
    \begin{cases}
        0, &\text{i = 0, j = 0}\\
        i, &\text{j = 0, i > 0}\\
        j, &\text{i = 0, j > 0}\\
        \min
        \begin{cases}
            D(i, j - 1) + 1,\\
            D(i - 1, j) + 1,\\
            D(i - 1, j - 1),\\
            D(i - 2, j - 2) + 1,
        \end{cases} 
        &\begin{aligned}
            & \text{если $i, j > 1$}, \\
            & S_{1}[i] = S_{2}[j - 1], \\
            & S_{1}[i - 1] = S_{2}[j],
        \end{aligned} \\
        \min
        \begin{cases}
            D(i - 1, j) + 1, \\
            D(i, j - 1) + 1, \\
            D(i - 1, j - 1) + \text{m}(S_1[i], S_2[j])
        \end{cases} &\text{иначе.}
    \end{cases}
    \label{eqn:recur-damlev}
\end{multline}

\section*{Вывод}
В данной части были рассмотрены формулы Левенштейна и Дамерау~---~Левенштейна, используемые для вычисления расстояния между строками. Обе формулы задаются рекуррентно, что позволяет реализовать как рекурсивный, так и итеративный подходы.